\begin{recipe}{Flammkuchen}
\kicker[Hoofdgerecht,Vegetarisch,Elzassisch]{Hoofdgerecht · vegetarisch · Elzassisch}
\index[register]{Flammkuchen}
\index[register]{Geitenkaas}
\index[register]{Crème fraîche}

\meta{25 minuten \dvd Oven 250\,°C of hoger}

\lettrine{F}{lammkuchen} is een dunne, knapperige bodem uit de Elzas, ook bekend als
tarte flambée. Met dit basisdeeg maak je twee flammkuchen, en de toppings
wissel je naar smaak. Dit is ook erg lekker om later koud te eten als hap bij de
borrel of de lunch.

\blockrule{Ingrediënten}

\noindent\textit{Deeg}
\begin{ingredients}
  \ing{300 g}{bloem}
  \ing{150 ml}{handwarm water}
  \ing{60 ml}{olijfolie}
  \ing{½ tl}{zout}
\end{ingredients}

\medskip
\noindent\textit{Traditionele topping}
\begin{ingredients}
  \ing{200 g}{crème fraîche}
  \ing{200 g}{spekreepjes}\index[register]{Spek}
  \ing{1}{rode ui, in ringen}
  \ingb{bieslook, fijngehakt}
\end{ingredients}

\medskip
\noindent\textit{Topping honing en geitenkaas}
\begin{ingredients}
  \ing{100 g}{geitenkaas, verkruimeld}\index[register]{Geitenkaas}
  \ing{50 g}{walnoten, grof gehakt}
  \ingb{honing}
\end{ingredients}

\blockrule{Bereiding}
\tip{Een hete pizzasteen geeft de knapperigste bodem. Flammkuchen is ook koud
lekker, dus restjes bewaar je gewoon.}
\begin{steps}
  \item Verwarm de oven zo heet mogelijk voor (250\,°C of hoger, boven- en
        onderwarmte). Leg een pizzasteen
        in de oven als je die hebt.
  \item Meng bloem, water, olijfolie en zout tot een soepel deeg. Kneed een paar
        minuten door.
  \item Laat het deeg 10 minuten rusten onder een vochtige theedoek. 
  \item Verdeel het deeg in 2 ballen. Rol elke bal zo dun mogelijk uit op een
        bebloemd oppervlak. Een vel bakpapier eronder helpt bij het verplaatsen naar de oven.
  \item Beleg de flammkuchen: bestrijk de bodem dun met crème fraîche of honing.
        Verdeel de rest van de topping erover.
  \item Schuif de flammkuchen op de hete steen of een bakplaat. Bak 8--10 minuten
        tot de randen knapperig zijn.
\end{steps}

\heroimagefade{images/flammkuchen}
\end{recipe}
